\chapterimage{figs/chapterhead.png}
\chapter{Applications, Tools, \texttt{C++} Examples, Tests, and Project Evolution}
\label{chap:applications_tools_exemplos_testes_evolucao}

\section{Introduction}
\label{sec:cap10_introducao}

This chapter wraps up the main structure of this developer manual by broadening the focus again. In the previous chapters, the attention fell on the central architecture of the system: organization of the source code, kernel, structure of model elements, \textit{plugins}, and parser. Now the goal is to examine how this infrastructure projects into concrete applications, helper tools, programmatic examples in \texttt{C++}, and testing mechanisms. In other words, this chapter is about the development ecosystem that forms around the \textit{GenESyS} core.

The importance of this closure is great. An extensible simulator doesn't just live off its core classes. It also depends on applications that expose it to the user, on tools that support its use and maintenance, on examples that show how the library can be used programmatically and on tests that support the safe evolution of the system. Without these elements, the core of the software could exist, but the project as a development platform would be impoverished.

Therefore, this chapter should be read as a final synthesis of the second half of the book. It shows how GenESyS transitions:
\begin{itemize}
 \item from the kernel to applications;
 \item of the architecture for programmatic use;
 \item functional extension for technical maintenance;
 \item of stable infrastructure for the future evolution of the project.
\end{itemize}

\section{Applications as Projections of the Kernel}
\label{sec:cap10_applications_como_projecoes_kernel}

The \texttt{source/applications} folder is the place where GenESyS stops appearing just as infrastructure and starts to take on concrete forms of interaction. This observation is essential for the developer. It shows that the project was not designed to exist just as an abstract library, but to support real applications at its core.

Architecturally, this is very valuable. By separating the kernel from the applications, the system gains:
\begin{itemize}
 \item modularity;
 \item core reuse on different interfaces;
 \item lower coupling between simulation logic and form of interaction;
 \item greater capacity for evolution.
\end{itemize}

In practice, this separation allows the GUI, Shell application, model-specific applications, and other future layers to exist as distinct projections of the same set of core services.

\subsection{The GUI as an Application on the Core}
\label{subsec:cap10_gui_como_aplicacao_sobre_nucleo}

In the first half of this book, the GUI was treated as the main way of using the simulator. Now, it starts to be seen in another way: as an application built on the GenESyS infrastructure. This shift in perspective is especially important for the developer.

The GUI is not the core. It is an application that:
\begin{itemize}
 \item consumes the model;
 \item consumes the semantics of components and data;
 \item connects to simulation engines;
 \item exposes these capabilities to the user through a visual interface.
\end{itemize}

This separation is architecturally correct and explains why the system can, in principle, support other applications in addition to the graphical interface.

\subsection{The Shell and Model-Specific Applications as Alternative Forms of Use}
\label{subsec:cap10_shell_modelspecific_forma_alternativa}

The current build tree of the project shows that the command-line applications are handled in a particularly structured way. The \texttt{CMakeLists.txt} of \texttt{source/applications/shell} defines the \texttt{genesys\_shell} target, while \texttt{source/applications/modelSpecific} defines \texttt{genesys\_modelspecific\_app} and lets the build select a concrete model-specific example file to compose the executable.

This is extremely important for the developer. It shows that the command-line path is not just an old auxiliary tool or a marginal way of use. It is an effective application of GenESyS and, more than that, a valuable environment for:
\begin{itemize}
 \item study the programmatic use of the simulator;
 \item build compact examples;
 \item develop lightweight applications based on the library;
 \item check behaviors without GUI mediation.
\end{itemize}

The \textit{Worker} service remains outside this application path as a local or private component only; this chapter does not provide public deployment instructions or imply that public exposure is approved.

\section{The examples in \texttt{C++} as developer material}
\label{sec:cap10_exemplos_cpp_material_desenvolvedor}

In the first half of the book, the examples in \texttt{C++} of the command-line applications were deliberately left out, because the user manual should not start with code. Now, however, they begin to occupy a central place. This is because these examples show one of the most important capabilities of GenESyS as software: its use as a simulation library.

The model-specific application's \texttt{CMakeLists.txt} explicitly shows the existence of build-selectable examples, including \textit{smart} examples and more elaborate teaching examples, which confirms the role of these files as study material and architectural demonstration.

\subsection{Why These Examples Are So Valuable}
\label{subsec:cap10_por_que_exemplos_tao_valiosos}

The examples in \texttt{C++} are valuable because they allow you to observe the simulator from another angle. In the GUI, the model appears as a visual and operational artifact. In the programmatic examples, it appears as an object explicitly constructed by the developer. This helps reveal:
\begin{itemize}
 \item how \textit{Simulator} is instantiated;
 \item how the model is built by code;
 \item how components and \textit{data definitions} are created and related;
 \item how execution is triggered;
 \item how the simulator library can be reused in new applications.
\end{itemize}

\subsection{The bridge between examples in \texttt{C++} and saved models}
\label{subsec:cap10_ponte_exemplos_cpp_modelos_salvos}

One of the most didactically rich features of the project is the possibility of relating:
\begin{itemize}
 \item examples built programmatically in \texttt{C++};
 \item models saved in the \texttt{models/} folder.
\end{itemize}

This bridge is extremely useful for the developer because it shows the same system from two perspectives:
\begin{itemize}
 \item as library object;
 \item as a modeling and usage artifact.
\end{itemize}

This correspondence helps to consolidate the understanding of GenESyS as a platform, and not just as a ready-made application.

\section{Tools as a Support Layer}
\label{sec:cap10_tools_como_camada_apoio}

The \texttt{source/tools} folder represents an important region of the project, although it is often less discussed than the kernel, parser, and \textit{plugins}. Its general function is to support system use, development, integration, or maintenance.

From a developer's point of view, tools are relevant for three reasons:
\begin{itemize}
 \item help make the simulator ecosystem more productive;
 \item offer points of support for observation or automation;
 \item can serve as a basis for new integrations.
\end{itemize}

In architectural terms, the existence of \texttt{tools} shows that the project is not limited to the ``kernel versus application'' dichotomy. There is also space for intermediate instruments, utilities and support mechanisms.

\subsection{When to study \texttt{tools}}
\label{subsec:cap10_quando_estudar_tools}

Tools should not always be the first reading priority. In general, it makes more sense to study them after the developer already understands:
\begin{itemize}
 \item the simulator core;
 \item the \textit{plugins} system;
 \item the parser;
 \item at least one concrete application, such as the GUI or the Shell.
\end{itemize}

After that, reading the tools tends to be much more productive, because the developer can better see their role in the project ecosystem.

\section{The Importance of Tests}
\label{sec:cap10_importancia_dos_testes}

In a project of the size and ambition of \textit{GenESyS}, the existence of tests is an essential part of the software's evolutionary architecture. This not only means that the project ``has tests'', but that the build system itself already treats them as a structured part of development.

\texttt{CMakeLists.txt} of \texttt{source/tests} explicitly shows the division between unit tests and \textit{smoke tests}, with an aggregated target \texttt{genesys\_tests} and the possibility of enabling or disabling the construction of smoke tests by build option

This organization is particularly valuable because it gives the developer a way to think about validation in layers:
\begin{itemize}
 \item unit tests for local behavior and minor invariants;
 \item \textit{smoke tests} for broader system health.
\end{itemize}

\subsection{Why Tests Are Part of the Architecture, Not Just the Routine}
\label{subsec:cap10_testes_parte_arquitetura}

It is common to think of testing only as an activity external to the main project. In GenESyS, however, they must be read in a stronger way: as part of the system's evolutionary infrastructure. This is particularly important because the software has:
\begin{itemize}
 \item relatively dense kernel;
 \item extension mechanisms;
 \item multiple applications;
 \item integration between parser, model, events and observability.
\end{itemize}

In a context like this, changing code without testing support tends to be risky. Therefore, the developer must see the \texttt{source/tests} folder as a direct ally in project maintenance.

\subsection{How the Developer Should Use the Tests}
\label{subsec:cap10_como_usar_testes}

A mature development practice in GenESyS must consider testing in at least three moments:
\begin{enumerate}
 \item before the change, to understand the current basis of behavior;
 \item during the change, to avoid local regressions;
 \item after the change, to validate integration and general sanity.
\end{enumerate}

This discipline is especially important when the change affects the kernel, parser, \textit{plugins} or applications.

\section{Build, Applications, and Behavior Selection}
\label{sec:cap10_build_aplicacoes_selecao_comportamentos}

Even though the build is not the guiding axis of the book, it continues to offer the developer an important set of signals about the project organization. The Shell and model-specific build files, for example, show that the system was designed to allow configurable selection of the behavior of the command-line executable, either by a fixed Shell entry point or by concrete examples under \texttt{source/applications/modelSpecific}.

This reveals something important about the project:
\begin{quote}
 GenESyS is not just a program; it is a platform upon which application behaviors can be chosen and assembled.
\end{quote}

This observation is of great interest to the developer because it reinforces the modularity of the system and helps to think about future applications about the core.

\section{Project Evolution and Contribution Paths}
\label{sec:cap10_evolucao_projeto_caminhos_contribuicao}

The combined reading of kernel, parser, \textit{plugins}, applications, tools and tests shows that GenESyS is a project in continuous evolution. This evolution can occur on several fronts:
\begin{itemize}
 \item kernel improvements;
 \item expansion of the set of components and \textit{data definitions};
 \item evolution of the language and parser;
 \item strengthening the GUI and other applications;
 \item expansion of auxiliary tools;
 \item consolidation of the test base.
\end{itemize}

From the developer's point of view, the most useful question is no longer:
\begin{quote}
 Where can I move?
\end{quote}
and becomes:
\begin{quote}
 At which layer of the system does my change make the most sense architecturally?
\end{quote}

This change in perspective is what distinguishes an opportunistic modification from a technically mature contribution.

\subsection{How to Choose an Entry Point for Contribution}
\label{subsec:cap10_como_escolher_ponto_entrada}

The choice of entry point depends largely on the type of contribution desired.

\begin{itemize}
 \item If the intention is to change the central dynamics of the simulation, the study must begin in the kernel.
 \item If the intention is to expand language and expressions, the study must begin with the parser.
 \item If the intention is to add modeling elements, reading should focus on \textit{plugins}, components and \textit{model data}.
 \item If the intention is to improve the user experience, the GUI and applications become central.
 \item If the intention is to increase robustness, testing and tracing become a priority.
\end{itemize}

This way of thinking greatly improves the quality of the contribution to the project.

\subsection{The Role of Examples and Tests in Evolution}
\label{subsec:cap10_papel_exemplos_testes_evolucao}

Examples and tests play complementary roles in software evolution.

\begin{itemize}
 \item Examples show \emph{como usar} the system and help you explore new capabilities.
 \item Tests help verify \emph{se o sistema continua correto} when being modified.
\end{itemize}

A mature GenESyS developer needs to value both. Without examples, the system loses its ability to demonstrate and learn. Without testing, you lose evolutionary security.

\section{An Integrated Final View of GenESyS}
\label{sec:cap10_visao_final_integrada_genesys}

Coming to the end of this version of the developer manual, it is useful to return to the integrated view of the project.

\textit{GenESyS} can be understood as:

\begin{itemize}
 \item a simulation kernel;
 \item a hierarchy of objects representing model, components and data;
 \item a language system and parser;
 \item an extensibility infrastructure by \textit{plugins};
 \item a set of applications built on top of the kernel;
 \item an ecosystem of examples, tools and tests that supports use and evolution.
\end{itemize}

This integrated view is important because it avoids too partial readings. GenESyS is not just GUI, not just parser, not just \textit{plugin framework}, not just simulation library. It is the articulation of all these dimensions.

\begin{table}[htbp]
 \centering
 \caption{Main Dimensions of GenESyS as a Software Platform.}
 \label{tab:cap10_dimensoes_principais_genesys}
 \small
 \begin{tabular}{|p{3.8cm}|p{7.4cm}|}
 \hline
 \textbf{Dimension} & \textbf{Role in the Project} \\
 \hline
 Kernel & Central Simulation Infrastructure \\
 \hline
 Parser & Language, expressions and textual semantics of the model \\
 \hline
 Plugins & Simulator extensibility and expansion of modeling elements \\
 \hline
 Applications & Concrete forms of use, such as GUI and Shell application \\
 \hline
 Tools & Support for use, integration and maintenance \\
 \hline
 Tests & Evolutionary security, validation and regression prevention \\
 \hline
 Examples & Material for demonstration, study and programmatic use of the library \\
 \hline
 \end{tabular}
\end{table}

\section{Final Considerations}
\label{sec:cap10_consideracoes_finais}

This chapter has closed the main structure of the developer's manual by showing that \textit{GenESyS} should not only be read for its central classes, but also for its ecosystem of applications, tools, examples and tests. The GUI now appears as an application on top of the core. The Shell application and the model-specific examples in \texttt{C++} reveal the use of the simulator as a library. Tools expand the development environment. The test tree provides the basis for safe evolution. And the general organization of the project shows that GenESyS should be thought of as a software platform, and not just as a ready-made simulator.

With this, the reader now has a sufficiently broad and technically grounded view to use, study and expand GenESyS in a mature way. From here, further development can follow different paths: review and improvement of the GUI, more detailed study of components and \textit{plugins}, evolution of the parser, strengthening of the testing infrastructure, development of new applications or creation of new modeling domains on the same core. In all these cases, the logic remains the same: understand the correct layer, respect the system architecture and evolve the simulator based on its own structural principles.

Test your understanding of this content by answering the following questions:

\begin{enumerate}
 \item Why should the GUI be understood, from the developer's point of view, as an application on the core and not as the entire system?

\item What is the architectural value of the examples in \texttt{C++} of the Shell and model-specific applications?

\item Why should the project's test base be treated as part of the software evolution infrastructure?

\item How does the distinction between kernel, parser, \textit{plugins}, applications, tools, examples and tests help the developer to better choose the entry point for a contribution?
\end{enumerate}
